\chapter{Agitation}

\section*{Definition}
Agitation is excessive motor or verbal activity with restlessness, irritability, distress, or threatening behavior. It is a symptom rather than a diagnosis. It becomes clinically urgent when it interferes with care, judgment, or safety. Agitation can occur with clear awareness, but new agitation with inattention or disorientation may be delirium from a medical cause.

\section*{When to See a Doctor}

\begin{itemize}
 \item Agitation is sudden, severe, rapidly worsening, or accompanied by confusion, reduced alertness, fever, severe headache, or altered mental status
 \item The person may harm themselves or someone else, cannot care for basic needs, or has access to weapons
 \item It follows a head injury, seizure, suspected overdose, or abrupt alcohol or sedative withdrawal
 \item It occurs with chest pain, severe shortness of breath, very high temperature, or signs of stroke such as facial drooping, one-sided weakness, or trouble speaking
\end{itemize}

\section*{How It Develops}
Agitation occurs when brain arousal, attention, emotion, or impulse control is disrupted. Pain, fear, sleep loss, intoxication or withdrawal, medication effects, and illness can all increase arousal or reduce the ability to regulate behavior. Delirium causes a fluctuating disturbance in attention and awareness; psychosis or mania may cause agitation while attention remains relatively intact. The same outward behavior can therefore have very different causes and treatments.

\section*{Causes}
\begin{itemize}
 \item \textbf{Psychiatric:} Anxiety, bipolar disorder with mania or a mixed episode, schizophrenia or another psychotic disorder, major depression, PTSD, and personality-related crises
 \item \textbf{Neurological:} Delirium, dementia, traumatic brain injury, stroke, epilepsy (during or after a seizure), Parkinson disease, and other brain disorders
 \item \textbf{Substance-related:} Alcohol or drug intoxication or withdrawal, stimulant toxicity, and occasionally excessive caffeine
 \item \textbf{Metabolic/medical:} Hyperthyroidism, hypoglycemia, hypoxia or high carbon dioxide, pain, infection, fever, dehydration, and electrolyte disturbances
 \item \textbf{Medication-induced:} Corticosteroids, antidepressants, anticholinergics, dopamine agonists, sedatives, and medication interactions
\end{itemize}

\section*{Related Symptoms}
\begin{itemize}
 \item Irritability, pacing, shouting, threatening behavior, insomnia, or inability to remain still
 \item Anxiety, panic, pressured speech, racing thoughts, hallucinations, or delusions
 \item Inattention, disorientation, fluctuating alertness, fever, tremor, sweating, or abnormal movements
\end{itemize}
\textbf{Important:} Agitation is a nonspecific symptom that requires identification of the underlying cause to guide management.

\section*{Medical Evaluation}
\begin{itemize}
 \item Immediate attention to safety, airway, breathing, circulation, vital signs, temperature, oxygen saturation, and bedside glucose, followed by a focused physical and neurologic examination
 \item History from the person and collateral sources about baseline behavior, onset, triggers, sleep, pain, recent illness, medications, missed doses, and alcohol or other substances
 \item Mental status examination assessing attention, orientation, psychosis, mood or mania, cognition, suicidal or homicidal thoughts, and capacity to participate in care
 \item Focused tests guided by the presentation, which may include blood count, electrolytes, kidney and liver function, glucose, calcium, thyroid testing, pregnancy testing when relevant, urinalysis, and toxicology testing. A toxicology screen does not identify every drug or replace clinical assessment.
 \item Neuroimaging for new focal neurologic signs, head trauma, seizure, severe headache, anticoagulant use, or unexplained delirium when intracranial disease is possible; EEG may be needed when non-convulsive seizure is suspected
\end{itemize}

\section*{Treatment Options}
\begin{itemize}
 \item Identify and treat the cause: correct hypoglycemia or hypoxia, treat pain or infection, manage delirium, and provide supervised treatment for alcohol or sedative withdrawal. Do not assume that agitation is psychiatric until urgent medical causes have been considered.
 \item Start with verbal de-escalation, a quiet low-stimulation setting, personal space, a calm introduction, simple choices, and attention to immediate needs. Use medication only when necessary for severe distress or imminent danger and when de-escalation is insufficient.
 \item In an emergency, a clinician may choose an oral, intramuscular, or intravenous medicine based on the likely cause and the person's medical risks. Options may include an antipsychotic, a benzodiazepine, or both; respiratory depression, low blood pressure, oversedation, drug interactions, and heart-rhythm effects require monitoring. Benzodiazepines are generally preferred for alcohol or sedative withdrawal and stimulant toxicity, but may worsen some intoxications or delirium.
 \item Physical restraint or seclusion should be a last resort, used only by trained staff under local law and policy, for the shortest time necessary with continuous monitoring of breathing and circulation.
 \item In dementia-related agitation, investigate pain, infection, constipation, medication effects, sleep, and environmental triggers first. Antipsychotics are reserved for severe agitation or psychosis that is dangerous or causes significant distress, at the lowest effective dose with regular review and a plan to stop.
\end{itemize}

\section*{Self-Care and Home Management}
\begin{itemize}
 \item If agitation is new, severe, or worsening, seek urgent medical help rather than trying to manage it alone. Keep a safe exit, remove weapons and hazardous objects, and do not physically restrain the person.
 \item Use a calm, non-confrontational tone, give the person space, offer simple choices, and avoid arguing about beliefs or blocking their movement unless immediate safety requires intervention.
 \item Reduce noise and crowding, maintain a familiar routine, and ask a trusted person to stay nearby if that is reassuring and safe.
 \item Avoid alcohol and recreational drugs. Do not abruptly stop alcohol, benzodiazepines, or other medicines that may cause withdrawal without medical advice. Do not encourage physical activity when the person is confused, medically unwell, or unsafe.
 \item If there are threats of self-harm or harm to others, call emergency services and stay nearby only if doing so is safe.
\end{itemize}

\section*{References}
\begin{thebibliography}{99}
\bibitem{agitation:ref1} Garriga M, Pacchiarotti I, et al. ``Assessment and management of agitation in psychiatry.'' \textit{J Clin Psychiatry}. 2016;77(6):e713--e724.
\bibitem{agitation:ref2} Zeller SL, Citrome L, et al. ``The psychopharmacology of agitation: consensus statement of the American Association for Emergency Psychiatry Project BETA.'' \textit{West J Emerg Med}. 2012;13(1):26--34.
\bibitem{agitation:ref3} American College of Emergency Physicians. Clinical policy: critical issues in the evaluation and management of adult out-of-hospital or emergency department patients presenting with severe agitation. Approved 2023.
\bibitem{agitation:ref4} Sadock BJ, Sadock VA, et al. \textit{Kaplan \& Sadock\'s Comprehensive Textbook of Psychiatry}, 10th ed. Wolters Kluwer, 2017.
\bibitem{agitation:ref5} National Institute for Health and Care Excellence. Violence and aggression: short-term management in mental health, health and community settings (NG10). Updated 2025.
\bibitem{agitation:ref6} American Psychiatric Association. The American Psychiatric Association practice guideline on the use of antipsychotics to treat agitation or psychosis in patients with dementia. 2016.
\end{thebibliography}
